% titel: Blatt 9, A9: N2+O2 zu 2NO, Gleichgewichtskonzentration
% skript: Mortimer Kap. 14 (Gleichgewicht, Hg: Memorie)

% Teil I: Werkzeuge, kein Lösungsweg.

Steht eine Gleichgewichtskonstante direkt zur Verfügung, löst man die
MWG-Gleichung nach der unbekannten Gleichgewichtskonzentration auf.

\begin{werkzeugkasten}[Gleichgewichtskonzentration berechnen]
Reaktion $N_2(g) + O_2(g) \rightleftharpoons 2NO(g)$.
Ausgang je $c_0$. Im Gleichgewicht:
\[
  [N_2]=[O_2] = c_0 - x,\qquad [NO]=2x .
\]
\[
  K_c = \frac{[NO]^2}{[N_2][O_2]}
      = \frac{(2x)^2}{(c_0-x)^2} .
\]
Nach $x$ auflösen (Wurzel oder quadratische Gleichung). $c_0 - x > 0$.
\textit{Skript:} Mortimer Kap. 14.
\end{werkzeugkasten}

\begin{achtung}[Darauf achten]
\begin{itemize}
  \item Hier entstehen 2 NO aus je 1 N$_2$ und O$_2$: die Änderung $x$ in
        mol/L wird von $N_2$/$O_2$ abgezogen, NO verdoppelt.
  \item Werte in mol/L; $K_c = 0.55$ (dimensionslos wg. gleicher
        Zähler-/Nennerzähl).
  \item Wurzel ziehen: $\frac{2x}{c_0-x} = \sqrt{K_c}$.
  \item Nach $x$ und damit $[NO]$ fragen, nicht nach einer Zwischengrösse.
\end{itemize}
\end{achtung}
